\section*{Chapitre \ref{ch:dv:fourier-pricing-and-calibration-engineering} --- Valorisation par transformée de Fourier et ingénierie de la calibration}

\begin{solution}{exo:dv:fourier-pricing-and-calibration-engineering:1}
$|\varphi(u-\frac i2)|\le\E[e^{x_T/2}]\le1$ par l'inégalité de Jensen, puisque $\E[e^{x_T}]=1$. L'intégrale est donc bornée par
$\int_0^\infty du/(u^2+\frac14)=\pi$, et $\sqrt{FK}\to0$ donne $C\to P(0,T)F$. $\varphi(-i)=\E[e^{x_T}]=\E[S_T]/F=1$ signifie que le forward est la
moyenne de $S_T$ sous la mesure de valorisation, comme il se doit pour un prix de type martingale forward.
\end{solution}

\begin{solution}{exo:dv:fourier-pricing-and-calibration-engineering:2}
Lorsque $k\to-\infty$, le call tend vers $F$, donc il n'est pas intégrable en $k$ et n'a pas de transformée de Fourier. Le facteur $e^{\alpha k}$ le fait décroître
de ce côté. La transformée fait alors intervenir $\varphi(v-(\alpha+1)i)$, c'est-à-dire $\E[(S_T/F)^{\alpha+1}]$, qui doit être fini~: $\alpha$ est limité par
les moments dont dispose le modèle. Dans le \omterm{def:dv:stochastic-volatility:heston}{modèle de Heston} avec une corrélation fortement négative et une longue échéance, les moments élevés explosent en temps fini.
\end{solution}

\begin{solution}{exo:dv:fourier-pricing-and-calibration-engineering:3}
Le prix du modèle est la formule de Black à sa volatilité implicite~: $C^{\text{modèle}}=C^{\text{Black}}(\sigma^{\text{modèle}})$. En développant autour de
$\sigma^{\text{milieu}}$, $C^{\text{modèle}}-C^{\text{milieu}}=\mathcal V(\sigma^{\text{modèle}}-\sigma^{\text{milieu}})+\frac12\,\partial_\sigma\mathcal
V\,(\sigma^{\text{modèle}}-\sigma^{\text{milieu}})^2+\dots$ En divisant par $\mathcal V$, il reste l'erreur de volatilité plus un terme du second ordre (le \omterm{def:dv:greeks-and-the-hedging-pnl:greeks}{volga}).
\end{solution}

\begin{solution}{exo:dv:fourier-pricing-and-calibration-engineering:4}
Avec $n$ cotations affectées d'un bruit indépendant $\sigma$ et $p$ paramètres ajustés, un modèle correct laisse des résidus de racine de moyenne quadratique
$\sigma\sqrt{(n-p)/n}$~: $0.3\sqrt{13/18}=0.255$ point de volatilité pour dix-huit cotations après le filtre. Les calibrations du chapitre donnent 0.22 le premier
jour et 0.291 en moyenne~: bruit auquel s'ajoute une légère erreur de spécification, puisque le marché est un \omterm{def:dv:jumps-and-levy-models:bates}{modèle de Bates}. Une calibration nettement inférieure au plancher du bruit reviendrait à
ajuster le bruit.
\end{solution}

\begin{solution}{exo:dv:fourier-pricing-and-calibration-engineering:5}
La méthode COS valorise l'option sous la densité tronquée à $[a,b]$. La masse située à l'extérieur, dans la queue gauche épaisse d'un \omterm{def:dv:stochastic-volatility:heston}{modèle de Heston} avec $\rho=-0.7$ et
$\eta=0.6$, est simplement perdue. Davantage de termes convergent vers le prix sous la densité tronquée, et non vers le vrai. Seul un intervalle plus large (seize
écarts-types) supprime l'erreur, au prix de davantage de termes pour une même résolution.
\end{solution}

\begin{solution}{exo:dv:fourier-pricing-and-calibration-engineering:6}
La \omterm{def:dv:stochastic-volatility:sv}{volatilité de la volatilité} se manifeste dans la courbure et le skew du smile de volatilité, amortie par le retour à la moyenne. Sur un horizon $T$, la réponse de la variance à
ses propres chocs est moyennée par le facteur $(1-e^{-\kappa T})/(\kappa T)$. Un $\eta$ plus grand associé à un $\kappa$ plus rapide donne à peu près la même
volatilité effective de la variance aux échéances cotées (jusqu'à un an). Les cotations ne voient que le produit, et non les deux séparément, d'où la vallée.
\end{solution}

\begin{solution}{exo:dv:fourier-pricing-and-calibration-engineering:7}
Avec des pondérations égales sur les prix, la plus forte variation quotidienne de $\eta$ est de 0.255 (0.235 avec une \omterm{def:dv:fourier-pricing-and-calibration-engineering:vega}{pondération par le vega}), et l'écart-type des variations quotidiennes de 0.066
(0.089). La dérive est à peu près la même. Les pondérations changent les cotations qui fixent les paramètres (avec des pondérations égales, les options à un an, dont les \omterm{def:dv:greeks-and-the-hedging-pnl:greeks}{vegas}
sont les plus élevés), et non l'aplatissement de la vallée. La pénalisation reste nécessaire.
\end{solution}

\begin{solution}{exo:dv:fourier-pricing-and-calibration-engineering:8}
L'erreur de calibration ne mesure que les vanilles. Dans le chapitre, les calibrations non pénalisées restent sous 0.45 point de volatilité chaque jour (0.42 au plus), tandis que
$\eta$ varie de 0.235 en un jour et l'ajustement de convexité d'un \omterm{def:dv:variance-swaps-and-volatility-derivatives:volswap}{swap de volatilité} de 0.93 point de volatilité. Il faut contrôler les trajectoires des paramètres, leurs alertes et
le \omterm{def:dv:fourier-pricing-and-calibration-engineering:recalib}{P\&L de recalibration} des exotiques aussi.
\end{solution}

\begin{solution}{pb:dv:fourier-pricing-and-calibration-engineering:1}
\textbf{1.} COS~: elle valorise directement les strikes cotés, et son erreur décroît exponentiellement avec le nombre de termes. La FFT donne une grille fixe et nécessite
une interpolation.
\textbf{2.} 128 termes ($8.6\times10^{-5}$).
\textbf{3.} Les cumulants~: $[a,b]=c_1+\ln(F/K)\pm L\sqrt{c_2}$ avec $L=16$. Trop étroit ($L=12$) et l'erreur stagne à $3.0\times10^{-5}$.
\textbf{4.} Environ huit itérations de six évaluations de résidus, chacune valorisant quatre smiles~: soit deux cents smiles, 800 évaluations de la fonction caractéristique
chacune. Ces évaluations (exponentielles complexes et logarithmes) dominent~; dans la construction, une calibration prend nettement moins d'une seconde.
\textbf{5.} Par comparaison avec la référence Gauss--Legendre de Lewis, avec la formule de Black pour une fonction caractéristique lognormale, et par $\varphi(0)=
\varphi(-i)=1$.
\textbf{6.} Deux cotations, l'une croisée et l'autre périmée (large)~; dix-huit subsistent.
\textbf{7.} $v_0=0.0451$, $\kappa=1.82$, $\bar v=0.0518$, $\eta=0.488$, $\rho=-0.632$~; 0.22 point de volatilité.
\textbf{8.} Son erreur passe de 0.35 à 0.46 point de volatilité.
\textbf{9.} Entre $\eta=0.35$ et 0.65, la meilleure erreur reste entre 0.23 et 0.29 point de volatilité. Sur toute la plage de 0.25 à 1, elle varie de
moins que le bruit des cotations, de 0.3 point.
\textbf{10.} $v_0$, $\rho$ et $\eta$ sont identifiés (écarts-types 0.0007, 0.046 et 0.049)~; $\kappa$ ne l'est pas (0.56, 28~\% de sa valeur).
\textbf{11.} 0.235, du jour 58 au jour 59 (de 0.51 à 0.28).
\textbf{12.} Elle baisse de 1.32 à 0.39 point de volatilité, soit une variation de 0.93~: 93\,000 pour un \omterm{def:dv:variance-swaps-and-volatility-derivatives:veganot}{notionnel vega} de 100\,000.
\textbf{13.} Les cotations ont à peine changé (l'erreur de calibration vaut 0.31 les deux jours)~; les paramètres ont glissé le long de la vallée. Aucun facteur de risque couvert par une couverture
n'a bougé.
\textbf{14.} La plus forte variation quotidienne de $\eta$ tombe à 0.071~; les variations quotidiennes de l'ajustement de convexité passent de 0.36 à 0.07 point de volatilité
(écart-type), et le jour 59 il varie de 0.18 au lieu de 0.93.
\textbf{15.} L'erreur de calibration moyenne passe de 0.291 à 0.303 point de volatilité, soit une hausse de 0.012.
\textbf{16.} Pour chaque $\lambda$ d'une grille allant de $10^{-5}$ à $10^{-2}$, rejouer les soixante jours avec le pipeline de production~; relever la plus forte variation quotidienne de
$\eta$ et l'erreur de calibration moyenne~; retenir le plus petit $\lambda$ qui respecte les deux limites.
\textbf{17.} Les paramètres retardent un véritable changement du marché (une hausse réelle de la \omterm{def:dv:stochastic-volatility:sv}{volatilité de la volatilité} après un choc). L'erreur de calibration augmente et le
modèle mal valorise les produits jusqu'à ce qu'il rattrape le retard. L'alerte sur l'erreur de calibration doit pouvoir passer outre la pénalisation.
\textbf{18.} Pénaliser davantage les directions mal identifiées ($\kappa$, $\eta$), et moins les directions bien identifiées ($v_0$, $\rho$), afin que celles-ci suivent
le marché chaque jour.
\textbf{19.} $\lambda=10^{-3}$~: la plus forte variation quotidienne de $\eta$ passe de 0.235 à 0.071 (à $3\times10^{-4}$, elle vaut encore 0.133), et l'erreur de calibration moyenne
augmente de 0.012 point de volatilité.
\textbf{20.} Des paramètres qui ne bougent que lorsque les cotations l'exigent, afin que les prix des exotiques suivent le marché et non le bruit.
\end{solution}

\begin{solution}{iq:dv:fourier-pricing-and-calibration-engineering:1}
Écrire le prix comme une intégrale dans l'espace de Fourier~: la \omterm{def:dv:fourier-pricing-and-calibration-engineering:lewis}{formule de Lewis} (identité de Parseval avec la transformée du payoff sur une bande), la \omterm{def:dv:fourier-pricing-and-calibration-engineering:cm}{formule de Carr--Madan}
transformée du call amorti évaluée par FFT, ou le développement COS de la densité. Toutes n'exigent que $\varphi$, et toutes sont vérifiées par $\varphi(0)=\varphi(-i)=1$
et par un cas lognormal comparé à Black.
\iqlookfor{l'une des trois voies, et une vérification.}
\end{solution}

\begin{solution}{iq:dv:fourier-pricing-and-calibration-engineering:2}
Carr--Madan~: une grille complète en log-strike en une seule FFT, mais un paramètre d'amortissement à choisir, une interpolation vers les strikes cotés, et des erreurs dues à l'espacement de la
grille et à son étendue. COS~: n'importe quels strikes, convergence exponentielle pour des densités régulières, peu de termes~; elle exige un intervalle de troncature issu des cumulants,
et les queues épaisses en demandent un large.
\iqlookfor{grille contre strikes, amortissement contre troncature, convergence.}
\end{solution}

\begin{solution}{iq:dv:fourier-pricing-and-calibration-engineering:3}
Le problème est mal posé~: une vallée de calibrations quasi équivalentes (pour Heston, $\kappa$ et $\eta$ ensemble), le long de laquelle le bruit déplace l'optimum. Le mesurer
(profil, recalibrations avec bruit, conditionnement de la jacobienne). Puis ajouter une pénalité sur les variations depuis la veille, choisie par rejeu, et figer ou pénaliser
les paramètres mal identifiés~; surveiller le \omterm{def:dv:fourier-pricing-and-calibration-engineering:recalib}{P\&L de recalibration} des exotiques.
\iqlookfor{caractère mal posé, pénalité, identification.}
\end{solution}

\begin{solution}{iq:dv:fourier-pricing-and-calibration-engineering:4}
Conserver l'instantané des cotations, les décisions du filtre, les réglages, le point de départ, le résultat et les diagnostics, afin que tout jour puisse être rejoué. Vérifier les
compteurs du filtre, l'erreur de calibration globale et par échéance, le conditionnement et les variations des paramètres. Déclencher une alarme dès qu'un paramètre dépasse sa limite ou
l'erreur de calibration au-delà de son propre seuil, et garder les paramètres précédents jusqu'à ce qu'une personne valide les nouveaux.
\iqlookfor{possibilité de rejeu, diagnostics, et une alarme qui bloque.}
\end{solution}

\begin{solution}{iq:dv:fourier-pricing-and-calibration-engineering:5}
Si les paramètres du modèle ont sauté le long d'une direction plate, les valeurs des exotiques ont bougé sans qu'aucun facteur de risque couvert ne varie~: \omterm{def:dv:fourier-pricing-and-calibration-engineering:recalib}{P\&L de recalibration}, que
l'attribution laisse sans explication. Vérifier l'historique des paramètres pour ce jour-là, revaloriser avec les paramètres de la veille sur le marché du jour, et comparer.
\iqlookfor{\omterm{def:dv:fourier-pricing-and-calibration-engineering:recalib}{P\&L de recalibration}, et revalorisation avec les paramètres de la veille.}
\end{solution}

\begin{solution}{iq:dv:fourier-pricing-and-calibration-engineering:6}
Des erreurs de prix non pondérées laissent dominer les options de longue échéance à la monnaie. Pondérer par l'inverse du \omterm{def:dv:greeks-and-the-hedging-pnl:greeks}{vega} pour calibrer en volatilité, et encore par l'inverse
de la fourchette (ou de la fiabilité des cotations), afin que les cotations liquides comptent davantage. Ajouter ensuite des pondérations pour ce dont le portefeuille a besoin, comme les échéances et les strikes des
couvertures de ses exotiques.
\iqlookfor{\omterm{def:dv:greeks-and-the-hedging-pnl:greeks}{vega}, fourchette, et les besoins du portefeuille.}
\end{solution}
